\caption{\textbf{a,} Three marginal trees of an illustrative ARG, observed alleles and reference population labels. Each sample has the same identity across trees. One SNP per interval is shown. The encoder pools sites within each interval. The cut $t_c=350$ generations defines groups of sampled descendants.
\textbf{b,} MLPs and site pooling initialize node features $b$. Each GNN application processes the complete tree $T_t$ upward then downward, with branch features setting message gates. Both applications share weights. Clade adapters describe structural changes and persistence. Their correction and memory correction $\delta$ modify $b$ before application (2).
\textbf{c,} Projected group states are copied to member haplotypes. A bidirectional exponential moving average (EMA) follows each haplotype along the genome, with decay determined by interval length and a learned distance scale. Group averaging, subtraction of the current projected group state (Center), and a linear map return $\delta$ to the group representative. Memory follows sampled haplotypes, not inferred ancestral identity.
\textbf{d,} Shown for one query $q$, a bidirectional GRU processes its local embeddings and interval geometry. Concatenating local and recurrent embeddings, an MLP and softmax give ancestry probabilities. Argmax produces the ancestry track. Colors and symbols are defined in the legend. Vectors and probabilities are illustrative, not fitted results.}
